%%% T04｜中值定理、洛必达与泰勒 %%%


\section*{2020 年}


\textbf{应用与证明题｜2 · 7 分}

已知 $f(x)$ 在 $[0,1]$ 上可导，且 $f(0) = f(1) = 0$，$f \left(\frac{1}{2}\right) = 2$，证明存在 $\xi \in (0,1)$，使 $f'(\xi) = 3$。


\section*{2021 年}


\textbf{计算题｜1 · 8 分}

求极限：$\lim_{x \to 0} \frac{2\cos x - 2 + x^2}{x(\sin x - x)}$


\textbf{证明题｜2 · 6 分}

设 $f(x)$ 在 $[1, 2]$ 上连续，在 $(1, 2)$ 内可导，且 $f(1) = f(2) = 0$，求证在 $(1, 2)$ 内至少存在一点 $\xi$，使得 $f'(\xi) + 3f(\xi) = 0$。


\section*{2022 年}


\textbf{计算题｜2 · 8 分}

设 $\xi$ 为函数 $f(x) = \arctan x$ 在区间 $[0, b]$ 上满足拉格朗日中值定理的“中值”，求极限 $\lim_{b \to 0} \frac{\xi}{b}$。


\textbf{证明题｜1 · 7 分}

设 $f(x)$ 在 $[1, 2]$ 上连续，在 $(1, 2)$ 内可导，且 $f(2) = 0$。证明：至少存在一点 $\xi \in (1, 2)$，使得 $\xi f'(\xi)\ln\xi + f(\xi) = 0$。


\section*{2023 年}


\textbf{选择题｜1 · 3 分}

设 $x \to 0$ 时，$x - \sin x$ 与 $x^n$ 为同阶无穷小，则 $n$ 等于（ \quad ）。\\
A. 1 \quad B. 2 \quad C. 3 \quad D. 4


\textbf{计算题｜1 · 8 分}

求极限 $\lim_{x \to 0} \left( \frac{\sin x}{x} \right)^{\frac{1}{x^2}}$。


\textbf{应用与证明题｜2 · 7 分}

设 $f(x)$ 在 $[0, 2]$ 上连续，在 $(0, 2)$ 内可导，且 $f(0)f(2) > 0$，$f(0)f(1) < 0$，试证至少存在一点 $\xi \in (0, 2)$，使 $f'(\xi) = f(\xi)$。


\section*{2024 年}


\textbf{计算题｜1 · 9 分}

已知 $\lim_{x \to 0}(x^{-3}\sin 3x + ax^{-2} + b) = 0$，求常数 $a$ 与 $b$。


\textbf{证明题｜1 · 13 分}

设函数 $f(x)$ 在 $\left[0, \frac{\pi}{4}\right]$ 上可导，且：$x \in \left[0, \frac{\pi}{4}\right]$ 时 $0 < f(x) < 1$，$x \in \left(0, \frac{\pi}{4}\right]$ 时 $f'(x) < 1$。证明：有且仅有一个 $x \in \left(0, \frac{\pi}{4}\right)$ 使得 $f(x) = \tan x$。


\section*{2025 年}


\textbf{选择题｜4 · 3 分}

函数 $f(x) = (2-x)^4$ 的三阶麦克劳林展式得余项 $R_3(x)$ 为（式中 $0 < \theta < 1$）（ \quad ）\\
A. $x^4$ \quad B. $\frac{1}{4!}(\theta x)^4$ \quad C. $\frac{-1}{4!}(\theta x)^4$ \quad D. $-x^4$


\textbf{计算题｜1 · 10 分}

当 $x \to 0$ 时，$ax^n$ 与 $\ln(1-x^3) + x^3$ 为等价无穷小，求常数 $a$ 与正整数 $n$。


\textbf{证明题｜2 · 10 分}

设 $f(x)$ 有二阶导数，且 $f(0) = f(1) = 0$，$\lim_{x \to 0} \frac{f(x)}{x} = 0$，则在 $(0, 1)$ 内至少存在一点 $\xi$，使 $f''(\xi) = 0$。
